% !TEX root = ../main.tex
\providecommand{\main}{..}
\documentclass[../main.tex]{subfiles}

\begin{document}
\chapter{Preliminary} \label{chap:preliminary}

\section{Compressibility of Deep Neural Networks}
In this section, we describe a simple way to quantify the compressibility of models. The source coding theorem~\cite{shannonInformationTheory} states that:
\begin{itemize}
    \item $n$ i.i.d. (independent and identically distributed) random variables, each with entropy $H(X)$, can be compressed to $nH(X)$ bits on average, with negligible information loss, as $n \rightarrow \infty$ 
    \item Conversely, if they are compressed to fewer bits on average, loss of information must occur.
\end{itemize}
In this case, $H(X)$, the entropy of a random variable, is the expected amount of information needed to describe it, and can be computed by the following.
\begin{align}
H(X)&:=-\sum_{x\in X}p(x)\log_2 p(x)
\end{align}
Note that the log base determines the units: $\log_2$ yields bits, $\log_{10}$ yields dits, and $\ln$ yields nats.

Consider a model with weights $W$, each stored in $b$ bits, and there are $n$ weights in total. The total number of bits occupied by the model is therefore $n\cdot b$. This is the capacity of the model: it is the maximum amount of information the weights could carry if every possible bit pattern were equally likely and the weights were statistically independent. 

In practice, however, models rarely use their full capacity. We can measure this by treating its weights as independent draws from a categorical distribution over these possible values.
\begin{align}
    \mathcal{X} = \{0, 1, \dots, 2^b-1\}
\end{align}

We can use Maximum Likelihood Estimation (MLE) to estimate the distribution as follows.
\begin{align}
    \hat{p}(x) &= \frac1n \sum^n_{k=1} \mathbf{1} (W_k=x) 
\text{, where } \mathbf{1}(p)=\begin{cases}
1 & \text{if }p\\
0 & \text{if }\neg p
\end{cases}
\end{align}
And we can estimate the entropy per weight as follows.
\begin{align}
    \hat{H}(W)&:=-\sum_{x\in \mathcal{X}}\hat{p}(x)\log_2 \hat{p}(x)
\end{align}

Note that if a certain value does not appear at all in the model weights, there will be an undefined term $0\ln 0$ in the above calculation. We use a convention in modern information theory to mitigate this, by using an extension of $x\ln x$ in the form of the following~\cite[eq. 1.1.2]{ebanks1998characterizations}.
\begin{align} \label{eq:math_hack}
    f(x) &= \begin{cases}
        0, & x = 0 \\
        x\ln x, & x > 0
    \end{cases}
\end{align}
This allows us to treat $0\ln 0$ as $0$. We can see that the total information stored in the weights would be approximately $n \cdot \hat{H}(W)$ bits. But the maximum possible total entropy of the model is $n\cdot b$ bits (achieved by the uniform distribution), it therefore follows that the capacity utilization by a model is the following.
\begin{align}
U(W) &:= \frac{\hat{H}(W)}{b}
\end{align}

Note that $U \in [0, 1]$, and $U = 1$ implies the model is nearly incompressible losslessly, while $U < 1$ implies the model can be losslessly compressed to $1 - U$ times its current size. For example, if $U = 0.5$, the model can be compressed to half its size with neglibile information loss.

Qwen 3~\cite{yangQwen3TechnicalReport2025b} is a state-of-the-art language model in various sizes. Now that we have a method to quantify the compressibility of models, let us evaluate the compressibility of the 4B model, which is provided in three data types, of widths 16, 8, and 4. We report the results in Table \ref{tab:qwen3_entropy}.
\todo{Consider reproducing with more models}

\begin{table}[H]
\centering
\caption{Entropy Estimate of Qwen 3 4B Weights}
\label{tab:qwen3_entropy}
\begin{tabular}{|c|c|c|}
 \hline
    Data Type & $\hat{H}(W)$ in bits & $U(W)$ \\ \hline
    bfloat16 & 10.58 & $10.58/16\approx66\%$ \\ \hline
    fp8 & 6.63 &  $6.63/8\approx82.9\%$ \\ \hline
    int4 & 3.67 & $3.67/4\approx91.8\%$ \\ \hline
\end{tabular}
\end{table}

We observe that the bfloat16 model, which is $\sim8$ GB on disk, only uses $\sim\frac23$ of its capacity. This implies that the model can be compressed to $\sim5.33$ GB with negligible information loss. This degree of compressibility motivates us to perform an investigation of the distribution of the weights by plotting their possible values against the probabilities of each value, which we report in Figures \ref{fig:qwen3_4b_bf16}, \ref{fig:qwen3_4b_fp8}, and \ref{fig:qwen3_4b_int4}.

\todo{No axis labels (need to rerun the code to add them)}

\begin{figure}[H]
    \centering
    \begin{subfigure}[b]{0.45\textwidth}
        \centering
        \includegraphics[width=\textwidth]{img/weight_distributions/qwen3_4b/bf16_exponent.png}
        \caption{Exponents}
    \end{subfigure}
    \hfill
    \begin{subfigure}[b]{0.45\textwidth}
        \centering
        \includegraphics[width=\textwidth]{img/weight_distributions/qwen3_4b/bf16_mantissa.png}
        \caption{Mantissas}
    \end{subfigure}
    \caption{Distribution of Qwen3-4B Weights (bfloat16)}
    \label{fig:qwen3_4b_bf16}
\end{figure}

\begin{figure}[H]
    \centering
    \begin{subfigure}[b]{0.45\textwidth}
        \centering
        \includegraphics[width=\textwidth]{img/weight_distributions/qwen3_4b/fp8_exponent.png}
        \caption{Exponents}
    \end{subfigure}
    \hfill
    \begin{subfigure}[b]{0.45\textwidth}
        \centering
        \includegraphics[width=\textwidth]{img/weight_distributions/qwen3_4b/fp8_mantissa.png}
        \caption{Mantissas}
    \end{subfigure}
    \caption{Distribution of Qwen3-4B Weights (fp8)}
    \label{fig:qwen3_4b_fp8}
\end{figure}

\begin{figure}[H]
\centering
\includegraphics[width=0.6\textwidth]{img/weight_distributions/qwen3_4b/int4.png}
\caption{Distribution of Qwen3-4B Weights (int4)}
\label{fig:qwen3_4b_int4}
\end{figure}

We notice that the full dynamic range of bfloat16 is actually not used. Especially in bfloat16, the exponents are clustered towards the center, with an average of 120 (this means most weights are $m \cdot 2^{-7}$). Several factors may contribute to this phenomenon: weight decay, regularization, and the optimizer (for example, stochastic gradient descent provides implicit optimization~\cite{smithOriginImplicitRegularization2021a}).

We therefore conclude that large language models often have much room for compression, and is not using all of its potential capacity. This provides motivation for research in compression methods.

\section{Small Language Models} \label{sec:slms}

In this section, we will examine TinyStories~\cite{eldanTinyStoriesHowSmall2023d}, a work that pushes the boundaries of small language models, and provide a few custom language models for further experimentation. By using a low-variance synthetic dataset (children's stories) generated by GPT-3.5 and GPT-4 containing only words a typical 3-4 year old would understand, they were able to generate coherent output using models with as few as 1 million parameters.

One limitation we observed in this work is their usage of the default GPT-Neo~\cite{blackGPTNeoLargeScale2021} tokenizer (which is identical to the pretrained GPT-2~\cite{radford2019language} tokenizer), restricted to the top 10,000 tokens. Using a generic tokenizer for a low-vocabulary dataset like TinyStories might result in wasted weights (tokens that never appear in TinyStories), and might not provide the optimal tokenzation for our specialized dataset. 

As a preliminary experiment, we train two small language models on TinyStories v1 and TinyStories v2, respectively. We use a SentencePiece Byte-Pair Encoding (BPE) tokenizer~\cite{kudoSentencePieceSimpleLanguage2018} with a vocabulary size of 1024, and we used a LLaMa-based~\cite{touvronLLaMAOpenEfficient2023c} model with roughly half the parameters of the smallest reference model.

\subsection{Evaluation: Quantitative Metric}
We use perplexity~\cite{jelinekPerplexityMeasureDifficulty1977} as a quantitative metric to evaluate a language model's accuracy. For a model $\theta$, a text $x$ of length $N$, it is defined as the following.
\begin{align}
    \text{Perplexity}(\theta, x) = \left( \prod_{i=1}^N \frac{1}{p_\theta(x_i\mid x_1, x_2, \dots, x_{i-1})} \right)^{\frac1N}
\end{align}
In other words, it is the geometric mean of the inverse probabilities of a given text. To show this with an example, suppose we have a model $\theta$ and a ground truth text ``Mary had a little lamb''. Suppose that the model's predicted probabilities of this text is the following.

\begin{table}[H]
\centering
\caption{Example Model Output Probabilities}
\begin{tabular}{|l|c|}
    \hline
    Item & Value \\ \hline
    $p_\theta(\text{``Mary''})$ & 0.5 \\ \hline
    $p_\theta(\text{``had''}\mid\text{``Mary''})$ & 0.7 \\ \hline
    $p_\theta(\text{``a''}\mid\text{``Mary'', ``had''})$ & 0.9 \\ \hline
    $p_\theta(\text{``little''}\mid\text{``Mary'', ``had'', ``a''})$ & 0.6 \\ \hline
    $p_\theta(\text{``lamb''}\mid\text{``Mary'', ``had'', ``a'', ``little''})$ & 0.9 \\ \hline
\end{tabular}
\end{table}

Then the perplexity of the model on the text is the following.
\begin{align}
    \text{Perplexity} &= \left(\frac{1}{0.5}\cdot\frac{1}{0.7}\cdot\frac{1}{0.9}\cdot\frac{1}{0.6}\cdot\frac{1}{0.9}\right)^{\frac15} \\
    &\approx 1.49
\end{align}

Intuitively, perplexity is a measure of how surprised or confused a model is when seeing a certain text. We compute the model's perplexity over its validation data, then use the result as a metric to evaluate the model's accuracy (lower is better).

\subsection{Evaluation: Results and Discussion}
We report the validation perplexity of our two models. Note that the original paper only provided models trained on TinyStories v1.

\begin{table}[H]
\centering
\caption{Evaluation of Our Language Models}
\label{tab:our_models_eval}
\begin{tabular}{|c|cc|c|}
\hline
                                       & \multicolumn{1}{c|}{Reference}        & Ours (TinyStories v1) & Ours (TinyStories v2) \\ \hline
Architecture                           & \multicolumn{1}{c|}{GPT-Neo}          & LLaMa                 & LLaMa                 \\ \hline
Parameter Count                          & \multicolumn{1}{c|}{1.17M}            & 0.57M                 & 0.46M                 \\ \hline
Training Data                                   & \multicolumn{2}{c|}{TinyStories v1}                  & TinyStories v2        \\ \hline
Tokenizer (Vocab Size)                 & \multicolumn{1}{c|}{GPT-Neo}          & SentencePiece                & SentencePiece                \\ \hline
Vocab Size                             & \multicolumn{1}{c|}{10000}            & 1024                  & 1024                  \\ \hline
Perplexity $\downarrow$ (TinyStories v1) & \multicolumn{1}{c|}{19.3919} & \textbf{11.9006}      & ---                   \\ \hline
Perplexity $\downarrow$ (TinyStories v2)& \multicolumn{1}{c|}{---}              & ---                   & 8.0994                \\ \hline
\end{tabular}
\end{table}

\todo{Add qualitative results (actual model outputs) here}

We observe that our model results in better perplexity than the reference model despite being much smaller. In \verb|float32|, each weight is 4 bytes, so our TinyStories v1 model with 0.57M parameters is expected to take only $\sim2.17$MiB of disk space (the actual disk space, including the model's configuration, is $\sim2.22$MiB). Models of this scale can be suitable for devices with extremely low compute and/or memory.

We will use our models as a base to conduct several more experiments.
\todo{Consider an ablation study to train the same GPT-Neo model, with only a different tokenizer}

\section{Baseline Round-to-Nearest Quantization} \label{sec:rtn}

The baseline for post-training quantization is symmetric uniform quantization using round-to-nearest (RTN) rounding~\cite{nagelWhitePaperNeural2021b}. A model can be quantized to $b$ bits using Algorithm \ref{alg:rtn_quant}, then dequantized using Algorithm \ref{alg:rtn_dequant}.

\begin{algorithm}[H]
\caption{Baseline RTN Quantization}\label{alg:rtn_quant}
\begin{algorithmic}[1]
\Require $W$ (unquantized weights), $b$ (target bit width)
\Ensure Quantized model $W_q$
\State $q_{\text{min}} \gets -2^{b - 1}$
\State $q_{\text{max}} \gets 2^{b - 1} - 1$
\State $\text{scale} \gets \frac{\max(W)}{q_{\text{max}}}$
\State $W_q \gets \text{clamp}\left(\lceil \frac{W}{\text{scale}} \rfloor, q_\text{min}, q_\text{max}\right)$
\State \Return $W_q$, scale
\end{algorithmic}
\end{algorithm}

\begin{algorithm}[H]
\caption{Baseline RTN Dequantization}\label{alg:rtn_dequant}
\begin{algorithmic}[1]
\Require $W_q$ (quantized weights), scale
\Ensure Dequantized weights $\hat{W}$
\State $\hat{W} \gets W_q \cdot \text{scale}$
\State \Return $\hat{W}$
\end{algorithmic}
\end{algorithm}

\section{Quantization is Weight Perturbation} \label{sec:quant_is_perturb}

Let us establish the observation that quantization results in a perturbed model, using an example. Consider a model $\theta$ with parameters $W = \begin{bmatrix}
1.11 \\ 2.22 \\ 3.33
\end{bmatrix}$. By applying the baseline round-to-nearest quantization from Section \ref{sec:rtn}, we obtain the following quantized model.
\begin{align}
    \text{scale} \approx 0.0262, W_{q} = \begin{bmatrix}
\texttt{0b00101010}\\ \texttt{0b01010101} \\ \texttt{0b01111111}
\end{bmatrix} = \begin{bmatrix}
42 \\ 85 \\ 127
\end{bmatrix}
\end{align}
After dequantization, we obtain the following model.
\begin{align}
W_{\text{reconstructed}} = \text{scale} \cdot W_{\text{quantized}} \approx \begin{bmatrix}
1.1013 \\ 2.2287 \\ 3.33
\end{bmatrix}
\end{align}
We can see that the model is perturbed by the following.
\begin{align}
    \Delta W = \begin{bmatrix}
1.1013 \\ 2.2287 \\ 3.33
\end{bmatrix} - \begin{bmatrix}
1.11 \\ 2.22 \\ 3.33
\end{bmatrix}
=\begin{bmatrix}
-0.0087 \\ 0.0087 \\ 0 \end{bmatrix}
\end{align}
It is important to note the direction of perturbation. In this case, our model moved in the direction of $\big[\begin{smallmatrix}
-1 \\ 1 \\ 0 \end{smallmatrix}\big]$.

\section{Effects of Perturbation of Model Weights}

As we established in Section \ref{sec:quant_is_perturb}, quantization perturbs weights in a certain direction. In this section, we examine the effects of weight perturbation using our TinyStories v2 model, by viewing its loss landscape.

\subsection{Language Model Optimization Process}
To understand our model's loss landscape, we must first understand its loss function. The goal in optimizing language models is to minimize the difference between the ground truth and the model's predictions. We measure this using the Kullback-Leibler (KL) divergence~\cite{kullbackInformationSufficiency1951a}, defined as follows.
\begin{align}
    D_{KL}(P\parallel Q)= \sum_{x\in X} p_P(x)\ln\frac{p_P(x)}{p_Q(x)}
\end{align}
The KL-divergence measures how an approximating distribution Q is different from a distribution P. We use one-hot encoding to encode the ground truth into the distribution P, and use the model's outputs as the approximating distribution Q. 

Note that we once again need to employ the convention in Equation \ref{eq:math_hack}. Let us show why with an example: if we are predicting a word ``Hello'', the ground truth $P$ distribution and the predicted $Q$ distribution might be:
\begin{table}[H]
\centering
\caption{Example of KL-divergence}
\begin{tabular}{|c|c|c|}
\hline
$\hat{y}$ & $p_P(\hat{y})$ & $p_Q(\hat{y})$ \\ \hline
``Today'' & 0 & 0.07 \\ \hline
``Hello'' & 1 & 0.81\\ \hline
``Cat'' & 0 & 0.10\\ \hline
``Rain'' & 0 & 0.02 \\ \hline
\end{tabular}
\end{table} 

In this case, the KL-divergence is:
\begin{align}
D_{KL}(P \parallel Q)
&= \underbrace{0 \cdot \ln\frac{0}{0.07}}_{\text{``Today''}}
+ \underbrace{1 \cdot \ln\frac{1}{0.81}}_{\text{``Hello''}}
+ \underbrace{0 \cdot \ln\frac{0}{0.10}}_{\text{``Cat''}}
+ \underbrace{0 \cdot \ln\frac{0}{0.02}}_{\text{``Rain''}},
\end{align}
which contains three undefined terms in the form of $0\cdot\ln(\frac{0}{n})$. 
% perhaps just 0\cdot \ln \frac0n = 0\ln 0
% but can I do \ln \frac0n = \ln 0?? idk
Using the convention from Equation \ref{eq:math_hack}, we can compute the undefined terms as follows.
\begin{align}
0\cdot \ln \frac{0}{n} = \underbrace{0\ln 0}_{0\text{ by convention}} - 0\ln n = 0
\end{align}

We then obtain the following.
\begin{align}
D_{KL}(P \parallel Q) &= -\ln(0.81) \approx 0.21
\end{align}

When minimizing the KL-divergence between the model predictions and the ground truth, we can use the following observation to ignore the constant $H(P)$ term, or the entropy of the ground truth distribution P.
\begin{align}
    D_{KL}(P\parallel Q)&= \underbrace{\sum_{x\in X} p_P(x)\ln p_P(x)}_{-H(P)} - \sum_{x\in X} p_P(x)\ln p_Q(x)\\
    &\propto -\sum_{x\in X} p_P(x)\ln p_Q(x)
\end{align}
When optimizing, we only need to minimize $-\sum_{x\in X} p(x)\ln q(x)$. This is also known as the cross entropy $H(P, Q)$ between distributions P and Q.

The cross entropy for a language model $\theta$ on a single token, where $P_i$ is the one-hot encoded ground truth and $Q_i$ is the model's predicted distribution, is the following.
\begin{align}
    H(P_i,Q_i) &= -\sum_{x \in X} p_{P_i}(x)\ln p_{Q_i}(x) = - \ln p_\theta (x_i \mid x_{<i})
\end{align}
The loss function, the average cross-entropy per token over a sequence of length N, is the following.
\begin{align}
    \mathcal{L}(\theta) = \frac1N \sum_{i=1}^N H(P_i, Q_i) = -\frac1N \sum_{i=1}^N \ln p_\theta (x_i \mid x_{<i})
\end{align}

Note that this is exactly the natural logarithm of the perplexity of the model on the given text. Given this loss function, we will proceed to investigate the loss landscape of our model.

\subsection{Loss Landscape}
Since the loss landscape of a language model is in $n$-dimension, where $n$ is its number of parameters, it is difficult to be directly visualized. We will investigate the loss landscape of our TinyStories v2 model from Section \ref{sec:slms} by perturbing its weights and computing the perplexity of the perturbed models.

\subsubsection{Straight Line}
First, we will sample models along a simple line in the following way.
\begin{align}
    \theta' = \theta + \epsilon,  \epsilon \in \{-0.0020 \dots 0.0015\} 
\end{align}
Note that the blue line in the figure represents $\epsilon \in \{-0.0015 \dots 0.0015\}$. From the plot in Figure \ref{fig:loss_line}, we notice that there is a ``basin'' of low loss, and that moving towards the negative direction caused less degradation than moving towards the positive direction.

\begin{figure}[H]
    \centering
    \includegraphics[width=0.6\textwidth]{img/loss_landscapes/loss_line.png}
    \caption{Models Along a Line}
    \label{fig:loss_line}
\end{figure}

\subsubsection{Moving Randomly}
Next, we sample models stochastically from the following distribution. We sample 5 models for each $\sigma$.
\begin{align}
    \theta' \sim \mathcal{N}(\theta, \sigma^2\mathbf{I}), \sigma^2\in\{0.0005\dots0.0040\}
\end{align}

Note that we cap the perplexity at 100 in Figure \ref{fig:loss_stochastic_capped} for better visualization, because some perplexities are extremely large.

\begin{figure}[H]
    \centering
    \includegraphics[width=0.6\textwidth]{img/loss_landscapes/loss_stochastic_capped.png}
    \caption{Models from Gaussian Distribution}
    \label{fig:loss_stochastic_capped}
\end{figure}

In Figure \ref{fig:loss_stochastic_avg}, we plot the average perplexities of each $\sigma$. We notice an outlier at $\sigma = 0.0024$, and removing it yields Figure \ref{fig:loss_stochastic_avg_outlier}.

\begin{figure}[H]
    \centering
    \begin{subfigure}[b]{0.45\textwidth}
        \centering
        \includegraphics[width=\textwidth]{img/loss_landscapes/loss_stochastic_avg.png}
        \caption{Averages}
        \label{fig:loss_stochastic_avg}
    \end{subfigure}
    \hfill
    \begin{subfigure}[b]{0.45\textwidth}
        \centering
        \includegraphics[width=\textwidth]{img/loss_landscapes/loss_stochastic_avg_outlier.png}
        \caption{Averages (Outlier Removed)}
        \label{fig:loss_stochastic_avg_outlier}
    \end{subfigure}
    \caption{Loss Landscape of Models from Gaussian Distribution}
\end{figure}

This further confirms the observation that when perturbing a model, the direction of perturbation is important, and some directions yield terrible results (like the outlier) while some directions preserve the model's accuracy more.

\subsubsection{Towards Another Model}
We pretrain another model with the same configuration as the base model, and move from the base model to the other model. The distance between these models, $\lVert \theta - \theta' \rVert_2$, is $75.6$, which means they're quite far apart, and the magnitude of perturbation is much greater than the previous experiments.

Note that when the x-axis is 0, the model is exactly the base model, and when the x-axis is 1, the model is exactly the other pretrained model.

\begin{figure}[H]
    \centering
    \includegraphics[width=0.6\textwidth]{img/loss_landscapes/loss_towards.png}
    \caption{Moving Towards Another Model}
    \label{fig:loss_towards}
\end{figure}

We observe from Figure \ref{fig:loss_towards} that even though the magnitude of perturbation is much greater than the previous experiments, the increase is loss is not as significant. This shows that the model is less sensitive to perturbations in this direction.

\todo{Consider adding an experiment: loss when moving towards the direction of the gradient}

\subsection{Discussion}
The key insight from this experiment is the fact that

\todo{Write how some directions are better than others, and worst directions cause significant degradation}

\end{document}
